% =====================================================================
% Esquema del proceso — TP2, Sección Cocinas de TODO DECO (punto a).
% Es la versión prolija del diagrama de docs/plan_de_trabajo.md §9.
%
% Uso en el informe:
%     \usepackage{tikz}
%     \usetikzlibrary{positioning, arrows.meta, calc, decorations.pathreplacing}
%     ...
%     \begin{figure}[htbp]
%       \centering
%       % =====================================================================
% Esquema del proceso — TP2, Sección Cocinas de TODO DECO (punto a).
% Es la versión prolija del diagrama de docs/plan_de_trabajo.md §9.
%
% Uso en el informe:
%     \usepackage{tikz}
%     \usetikzlibrary{positioning, arrows.meta, calc, decorations.pathreplacing}
%     ...
%     \begin{figure}[htbp]
%       \centering
%       % =====================================================================
% Esquema del proceso — TP2, Sección Cocinas de TODO DECO (punto a).
% Es la versión prolija del diagrama de docs/plan_de_trabajo.md §9.
%
% Uso en el informe:
%     \usepackage{tikz}
%     \usetikzlibrary{positioning, arrows.meta, calc, decorations.pathreplacing}
%     ...
%     \begin{figure}[htbp]
%       \centering
%       % =====================================================================
% Esquema del proceso — TP2, Sección Cocinas de TODO DECO (punto a).
% Es la versión prolija del diagrama de docs/plan_de_trabajo.md §9.
%
% Uso en el informe:
%     \usepackage{tikz}
%     \usetikzlibrary{positioning, arrows.meta, calc, decorations.pathreplacing}
%     ...
%     \begin{figure}[htbp]
%       \centering
%       \input{../docs/esquema_proceso.tex}
%       \caption{Del stock del depósito a la variedad de combos ofrecidos.}
%     \end{figure}
%
% Los colores son los mismos de los gráficos (src/graficos.py).
%
% Nota: el esquema no usa los caracteres "<" ni ">" en ningún lado (las flechas se
% declaran como -{Stealth} y las desigualdades con \leq). Es a propósito: babel en
% español los toma como atajos de «comillas» y rompería la compilación.
% =====================================================================

\definecolor{tpAcento}{HTML}{2A78D6}
\definecolor{tpAcentoSuave}{HTML}{CDE2FB}
\definecolor{tpGris}{HTML}{898781}
\definecolor{tpGrisSuave}{HTML}{E1E0D9}
\definecolor{tpTinta}{HTML}{0B0B0B}
\definecolor{tpTintaDos}{HTML}{52514E}

\begin{tikzpicture}[
    font=\footnotesize,
    espacio/.style={
      draw=tpGrisSuave, fill=tpGrisSuave!35, rounded corners=2pt,
      text=tpTinta, minimum width=4.9cm, minimum height=0.48cm,
      inner sep=3pt, align=center,
    },
    limita/.style={espacio, draw=tpAcento, fill=tpAcentoSuave},
    caja/.style={
      draw=tpGris, rounded corners=3pt, text=tpTinta,
      minimum width=3.1cm, minimum height=1.5cm, inner sep=5pt, align=center,
    },
    resultado/.style={caja, draw=tpAcento, fill=tpAcentoSuave},
    flecha/.style={-{Stealth[length=4pt]}, draw=tpGris, line width=0.7pt},
    etiqueta/.style={text=tpTintaDos, font=\scriptsize, midway, above},
  ]

  % --- Columna 1: los espacios del depósito (recurso) ---------------
  \node[espacio] (e1) {Baldosas \quad $\leq 5$};
  \node[espacio, below=1.6mm of e1] (e2) {Empapelado \quad $\leq 8$};
  \node[espacio, below=1.6mm of e2] (e3) {Apliques de luz \quad $\leq 4$};
  \node[espacio, below=1.6mm of e3] (e4) {Alacenas \quad $\leq 4$};
  \node[limita,  below=1.6mm of e4] (e5) {\textbf{Mesadas} \quad $\boldsymbol{\leq 3}$};
  \node[espacio, below=1.6mm of e5] (e6) {Bacha y grifería \quad $\leq 4$};
  \node[espacio, below=1.6mm of e6] (e7) {Lavavajillas + cocinas \quad $\leq 5$};

  \node[text=tpTintaDos, font=\scriptsize\bfseries, above=2.5mm of e1] (tituloDep)
    {DEPÓSITO (capacidad por espacio)};

  % Llave que agrupa los siete espacios.
  \coordinate (depTop) at ($(e1.north east)+(2mm,0)$);
  \coordinate (depBot) at ($(e7.south east)+(2mm,0)$);
  \draw[draw=tpGris, decorate, decoration={brace, amplitude=4pt}]
    (depTop) -- (depBot);

  % --- Columna 2: el stock (variables de decisión) ------------------
  \node[caja, right=15mm of e4] (stock)
    {\textbf{Stock}\\[1mm] $x_p \in \mathbb{Z}_{\geq 0}$\\[0.5mm] unidades de cada\\ uno de los 30 artículos};

  % --- Columna 3: los combos ----------------------------------------
  \node[caja, right=15mm of stock] (combos)
    {\textbf{Combos completos}\\[1mm] $y_c \in \{0,1\}$\\[0.5mm] 20 combos\\ de 7 u 8 artículos};

  % --- Columna 4: el resultado --------------------------------------
  \node[resultado, right=15mm of combos] (variedad)
    {\textbf{Variedad}\\[1mm] $\max \sum_c y_c$\\[0.5mm] combos disponibles\\ para la venta};

  % --- Flechas -------------------------------------------------------
  \draw[flecha] ($(depTop)!0.5!(depBot)+(4mm,0)$) -- (stock.west)
    node[etiqueta, align=center] {R2\\ capacidad};
  \draw[flecha] (stock.east) -- (combos.west)
    node[etiqueta, align=center] {R1\\ cobertura};
  \draw[flecha] (combos.east) -- (variedad.west);

  % --- Nota al pie del esquema ---------------------------------------
  \node[text=tpTintaDos, font=\scriptsize, align=center, below=6mm of combos] (nota)
    {Un combo cuenta solo si \emph{todos} sus artículos tienen\\
     una unidad reservada para él (supuesto S1).};

\end{tikzpicture}

%       \caption{Del stock del depósito a la variedad de combos ofrecidos.}
%     \end{figure}
%
% Los colores son los mismos de los gráficos (src/graficos.py).
%
% Nota: el esquema no usa los caracteres "<" ni ">" en ningún lado (las flechas se
% declaran como -{Stealth} y las desigualdades con \leq). Es a propósito: babel en
% español los toma como atajos de «comillas» y rompería la compilación.
% =====================================================================

\definecolor{tpAcento}{HTML}{2A78D6}
\definecolor{tpAcentoSuave}{HTML}{CDE2FB}
\definecolor{tpGris}{HTML}{898781}
\definecolor{tpGrisSuave}{HTML}{E1E0D9}
\definecolor{tpTinta}{HTML}{0B0B0B}
\definecolor{tpTintaDos}{HTML}{52514E}

\begin{tikzpicture}[
    font=\footnotesize,
    espacio/.style={
      draw=tpGrisSuave, fill=tpGrisSuave!35, rounded corners=2pt,
      text=tpTinta, minimum width=4.9cm, minimum height=0.48cm,
      inner sep=3pt, align=center,
    },
    limita/.style={espacio, draw=tpAcento, fill=tpAcentoSuave},
    caja/.style={
      draw=tpGris, rounded corners=3pt, text=tpTinta,
      minimum width=3.1cm, minimum height=1.5cm, inner sep=5pt, align=center,
    },
    resultado/.style={caja, draw=tpAcento, fill=tpAcentoSuave},
    flecha/.style={-{Stealth[length=4pt]}, draw=tpGris, line width=0.7pt},
    etiqueta/.style={text=tpTintaDos, font=\scriptsize, midway, above},
  ]

  % --- Columna 1: los espacios del depósito (recurso) ---------------
  \node[espacio] (e1) {Baldosas \quad $\leq 5$};
  \node[espacio, below=1.6mm of e1] (e2) {Empapelado \quad $\leq 8$};
  \node[espacio, below=1.6mm of e2] (e3) {Apliques de luz \quad $\leq 4$};
  \node[espacio, below=1.6mm of e3] (e4) {Alacenas \quad $\leq 4$};
  \node[limita,  below=1.6mm of e4] (e5) {\textbf{Mesadas} \quad $\boldsymbol{\leq 3}$};
  \node[espacio, below=1.6mm of e5] (e6) {Bacha y grifería \quad $\leq 4$};
  \node[espacio, below=1.6mm of e6] (e7) {Lavavajillas + cocinas \quad $\leq 5$};

  \node[text=tpTintaDos, font=\scriptsize\bfseries, above=2.5mm of e1] (tituloDep)
    {DEPÓSITO (capacidad por espacio)};

  % Llave que agrupa los siete espacios.
  \coordinate (depTop) at ($(e1.north east)+(2mm,0)$);
  \coordinate (depBot) at ($(e7.south east)+(2mm,0)$);
  \draw[draw=tpGris, decorate, decoration={brace, amplitude=4pt}]
    (depTop) -- (depBot);

  % --- Columna 2: el stock (variables de decisión) ------------------
  \node[caja, right=15mm of e4] (stock)
    {\textbf{Stock}\\[1mm] $x_p \in \mathbb{Z}_{\geq 0}$\\[0.5mm] unidades de cada\\ uno de los 30 artículos};

  % --- Columna 3: los combos ----------------------------------------
  \node[caja, right=15mm of stock] (combos)
    {\textbf{Combos completos}\\[1mm] $y_c \in \{0,1\}$\\[0.5mm] 20 combos\\ de 7 u 8 artículos};

  % --- Columna 4: el resultado --------------------------------------
  \node[resultado, right=15mm of combos] (variedad)
    {\textbf{Variedad}\\[1mm] $\max \sum_c y_c$\\[0.5mm] combos disponibles\\ para la venta};

  % --- Flechas -------------------------------------------------------
  \draw[flecha] ($(depTop)!0.5!(depBot)+(4mm,0)$) -- (stock.west)
    node[etiqueta, align=center] {R2\\ capacidad};
  \draw[flecha] (stock.east) -- (combos.west)
    node[etiqueta, align=center] {R1\\ cobertura};
  \draw[flecha] (combos.east) -- (variedad.west);

  % --- Nota al pie del esquema ---------------------------------------
  \node[text=tpTintaDos, font=\scriptsize, align=center, below=6mm of combos] (nota)
    {Un combo cuenta solo si \emph{todos} sus artículos tienen\\
     una unidad reservada para él (supuesto S1).};

\end{tikzpicture}

%       \caption{Del stock del depósito a la variedad de combos ofrecidos.}
%     \end{figure}
%
% Los colores son los mismos de los gráficos (src/graficos.py).
%
% Nota: el esquema no usa los caracteres "<" ni ">" en ningún lado (las flechas se
% declaran como -{Stealth} y las desigualdades con \leq). Es a propósito: babel en
% español los toma como atajos de «comillas» y rompería la compilación.
% =====================================================================

\definecolor{tpAcento}{HTML}{2A78D6}
\definecolor{tpAcentoSuave}{HTML}{CDE2FB}
\definecolor{tpGris}{HTML}{898781}
\definecolor{tpGrisSuave}{HTML}{E1E0D9}
\definecolor{tpTinta}{HTML}{0B0B0B}
\definecolor{tpTintaDos}{HTML}{52514E}

\begin{tikzpicture}[
    font=\footnotesize,
    espacio/.style={
      draw=tpGrisSuave, fill=tpGrisSuave!35, rounded corners=2pt,
      text=tpTinta, minimum width=4.9cm, minimum height=0.48cm,
      inner sep=3pt, align=center,
    },
    limita/.style={espacio, draw=tpAcento, fill=tpAcentoSuave},
    caja/.style={
      draw=tpGris, rounded corners=3pt, text=tpTinta,
      minimum width=3.1cm, minimum height=1.5cm, inner sep=5pt, align=center,
    },
    resultado/.style={caja, draw=tpAcento, fill=tpAcentoSuave},
    flecha/.style={-{Stealth[length=4pt]}, draw=tpGris, line width=0.7pt},
    etiqueta/.style={text=tpTintaDos, font=\scriptsize, midway, above},
  ]

  % --- Columna 1: los espacios del depósito (recurso) ---------------
  \node[espacio] (e1) {Baldosas \quad $\leq 5$};
  \node[espacio, below=1.6mm of e1] (e2) {Empapelado \quad $\leq 8$};
  \node[espacio, below=1.6mm of e2] (e3) {Apliques de luz \quad $\leq 4$};
  \node[espacio, below=1.6mm of e3] (e4) {Alacenas \quad $\leq 4$};
  \node[limita,  below=1.6mm of e4] (e5) {\textbf{Mesadas} \quad $\boldsymbol{\leq 3}$};
  \node[espacio, below=1.6mm of e5] (e6) {Bacha y grifería \quad $\leq 4$};
  \node[espacio, below=1.6mm of e6] (e7) {Lavavajillas + cocinas \quad $\leq 5$};

  \node[text=tpTintaDos, font=\scriptsize\bfseries, above=2.5mm of e1] (tituloDep)
    {DEPÓSITO (capacidad por espacio)};

  % Llave que agrupa los siete espacios.
  \coordinate (depTop) at ($(e1.north east)+(2mm,0)$);
  \coordinate (depBot) at ($(e7.south east)+(2mm,0)$);
  \draw[draw=tpGris, decorate, decoration={brace, amplitude=4pt}]
    (depTop) -- (depBot);

  % --- Columna 2: el stock (variables de decisión) ------------------
  \node[caja, right=15mm of e4] (stock)
    {\textbf{Stock}\\[1mm] $x_p \in \mathbb{Z}_{\geq 0}$\\[0.5mm] unidades de cada\\ uno de los 30 artículos};

  % --- Columna 3: los combos ----------------------------------------
  \node[caja, right=15mm of stock] (combos)
    {\textbf{Combos completos}\\[1mm] $y_c \in \{0,1\}$\\[0.5mm] 20 combos\\ de 7 u 8 artículos};

  % --- Columna 4: el resultado --------------------------------------
  \node[resultado, right=15mm of combos] (variedad)
    {\textbf{Variedad}\\[1mm] $\max \sum_c y_c$\\[0.5mm] combos disponibles\\ para la venta};

  % --- Flechas -------------------------------------------------------
  \draw[flecha] ($(depTop)!0.5!(depBot)+(4mm,0)$) -- (stock.west)
    node[etiqueta, align=center] {R2\\ capacidad};
  \draw[flecha] (stock.east) -- (combos.west)
    node[etiqueta, align=center] {R1\\ cobertura};
  \draw[flecha] (combos.east) -- (variedad.west);

  % --- Nota al pie del esquema ---------------------------------------
  \node[text=tpTintaDos, font=\scriptsize, align=center, below=6mm of combos] (nota)
    {Un combo cuenta solo si \emph{todos} sus artículos tienen\\
     una unidad reservada para él (supuesto S1).};

\end{tikzpicture}

%       \caption{Del stock del depósito a la variedad de combos ofrecidos.}
%     \end{figure}
%
% Los colores son los mismos de los gráficos (src/graficos.py).
%
% Nota: el esquema no usa los caracteres "<" ni ">" en ningún lado (las flechas se
% declaran como -{Stealth} y las desigualdades con \leq). Es a propósito: babel en
% español los toma como atajos de «comillas» y rompería la compilación.
% =====================================================================

\definecolor{tpAcento}{HTML}{2A78D6}
\definecolor{tpAcentoSuave}{HTML}{CDE2FB}
\definecolor{tpGris}{HTML}{898781}
\definecolor{tpGrisSuave}{HTML}{E1E0D9}
\definecolor{tpTinta}{HTML}{0B0B0B}
\definecolor{tpTintaDos}{HTML}{52514E}

\begin{tikzpicture}[
    font=\footnotesize,
    espacio/.style={
      draw=tpGrisSuave, fill=tpGrisSuave!35, rounded corners=2pt,
      text=tpTinta, minimum width=4.9cm, minimum height=0.48cm,
      inner sep=3pt, align=center,
    },
    limita/.style={espacio, draw=tpAcento, fill=tpAcentoSuave},
    caja/.style={
      draw=tpGris, rounded corners=3pt, text=tpTinta,
      minimum width=3.1cm, minimum height=1.5cm, inner sep=5pt, align=center,
    },
    resultado/.style={caja, draw=tpAcento, fill=tpAcentoSuave},
    flecha/.style={-{Stealth[length=4pt]}, draw=tpGris, line width=0.7pt},
    etiqueta/.style={text=tpTintaDos, font=\scriptsize, midway, above},
  ]

  % --- Columna 1: los espacios del depósito (recurso) ---------------
  \node[espacio] (e1) {Baldosas \quad $\leq 5$};
  \node[espacio, below=1.6mm of e1] (e2) {Empapelado \quad $\leq 8$};
  \node[espacio, below=1.6mm of e2] (e3) {Apliques de luz \quad $\leq 4$};
  \node[espacio, below=1.6mm of e3] (e4) {Alacenas \quad $\leq 4$};
  \node[limita,  below=1.6mm of e4] (e5) {\textbf{Mesadas} \quad $\boldsymbol{\leq 3}$};
  \node[espacio, below=1.6mm of e5] (e6) {Bacha y grifería \quad $\leq 4$};
  \node[espacio, below=1.6mm of e6] (e7) {Lavavajillas + cocinas \quad $\leq 5$};

  \node[text=tpTintaDos, font=\scriptsize\bfseries, above=2.5mm of e1] (tituloDep)
    {DEPÓSITO (capacidad por espacio)};

  % Llave que agrupa los siete espacios.
  \coordinate (depTop) at ($(e1.north east)+(2mm,0)$);
  \coordinate (depBot) at ($(e7.south east)+(2mm,0)$);
  \draw[draw=tpGris, decorate, decoration={brace, amplitude=4pt}]
    (depTop) -- (depBot);

  % --- Columna 2: el stock (variables de decisión) ------------------
  \node[caja, right=15mm of e4] (stock)
    {\textbf{Stock}\\[1mm] $x_p \in \mathbb{Z}_{\geq 0}$\\[0.5mm] unidades de cada\\ uno de los 30 artículos};

  % --- Columna 3: los combos ----------------------------------------
  \node[caja, right=15mm of stock] (combos)
    {\textbf{Combos completos}\\[1mm] $y_c \in \{0,1\}$\\[0.5mm] 20 combos\\ de 7 u 8 artículos};

  % --- Columna 4: el resultado --------------------------------------
  \node[resultado, right=15mm of combos] (variedad)
    {\textbf{Variedad}\\[1mm] $\max \sum_c y_c$\\[0.5mm] combos disponibles\\ para la venta};

  % --- Flechas -------------------------------------------------------
  \draw[flecha] ($(depTop)!0.5!(depBot)+(4mm,0)$) -- (stock.west)
    node[etiqueta, align=center] {R2\\ capacidad};
  \draw[flecha] (stock.east) -- (combos.west)
    node[etiqueta, align=center] {R1\\ cobertura};
  \draw[flecha] (combos.east) -- (variedad.west);

  % --- Nota al pie del esquema ---------------------------------------
  \node[text=tpTintaDos, font=\scriptsize, align=center, below=6mm of combos] (nota)
    {Un combo cuenta solo si \emph{todos} sus artículos tienen\\
     una unidad reservada para él (supuesto S1).};

\end{tikzpicture}
